%% kapa: exam=1 round=34 year=2023 period=2 form=A subject=appraisal_law subject_ko=감정평가관계법규 seq=9 number=9
%% source: Q-Net 감정평가사 기출문제 (cst00309) artlSeq=5211789 file=raw/1008/5211789/2023년도 제34회 감정평가사 1차시험 2교시.pdf
%% conversion: kordoc 4.24.51 -> tools/gen1.py (auto); review=proofread-sonnet; flags=none
\begin{question}{9}
국토의 계획 및 이용에 관한 법령상 도시·군관리계획의 결정권자가 입지규제최소구역으로 지정할 수 있는 지역에 해당하지 않는 것은?
\begin{choices}
  \item 세 개 이상의 노선이 교차하는 대중교통 결절지로부터 1킬로미터 이내에 위치한 지역
  \item 「도시재생 활성화 및 지원에 관한 특별법」에 따른 도시재생활성화지역 중 도시경제기반형 활성화계획을 수립하는 지역
  \item 「도시 및 주거환경정비법」에 따른 노후·불량건축물이 밀집한 상업지역으로 정비가 시급한 지역
  \item 창의적인 지역개발이 필요한 지역으로 「빈집 및 소규모주택 정비에 관한 특례법」에 따른 소규모주택정비사업의 시행구역
  \item 도시·군기본계획에 따른 도심·부도심
\end{choices}
\end{question}
\answer{3}
